\section{Evaluation Methodology}
\label{sec:evaluation-methodology}

\subsection{Evaluation Scope}
\label{subsec:evaluation-scope}

The evaluation of GEAR is organized around three complementary dimensions:
representational coverage and transformation fidelity, execution correctness and
behavioral fidelity, and usability and practical utility. The present evaluation
focuses on the first two dimensions. The third dimension will be investigated
separately through a participant-based study.

The first dimension evaluates the ability of GEAR to represent and preserve the
main design and orchestration capabilities of heterogeneous frameworks for
\gls{llm}-based \gls{mas}s. This dimension is examined in two directions. The
framework-to-GEAR direction determines whether the capabilities identified in the
ten selected frameworks can be expressed through the common variability model,
composed of \texttt{GearAgent}, \texttt{GearModule}, and
\texttt{GearWorkflow}. Conversely, the GEAR-to-framework direction determines
whether the properties expressed in a GEAR specification are preserved when they
are transformed into framework-specific implementations. This analysis considers
the mappings, transformations, and generated artefacts associated with each of the
ten frameworks presented in Table~\ref{tab:framework-metadata}.

The second dimension evaluates the generated artefacts at both the technical and
behavioral levels. Technical correctness concerns whether the generated artefacts
are syntactically valid, importable, and executable in their target frameworks.
Behavioral fidelity concerns whether their execution respects the orchestration
constraints expressed in the source GEAR specification, such as task ordering,
dependencies, parallel branches, synchronization points, conditional transitions,
and bounded loops. Controlled agent responses are used to isolate orchestration
behavior from the non-determinism of the underlying \gls{llm}s.

Cross-framework portability is examined across these two dimensions. The same
GEAR specifications are used as inputs for the different framework targets without
introducing framework-specific modifications into the source configurations. This
makes it possible to determine which properties are preserved across frameworks,
which require framework-specific adaptations, and which cannot be supported by a
given target. Unsupported or inapplicable capabilities remain part of the
evaluation and are reported explicitly rather than being excluded from the
analysis.

Table~\ref{tab:evaluation-scope} summarizes the dimensions considered, their
evaluation subjects, and the main forms of evidence used.

\begin{table*}[t]
    \centering
    \caption{Scope of the GEAR evaluation.}
    \label{tab:evaluation-scope}
    \footnotesize
    \begin{tabularx}{\textwidth}{
        @{}
        p{0.20\textwidth}
        p{0.26\textwidth}
        p{0.19\textwidth}
        >{\raggedright\arraybackslash}X
        @{}
    }
        \toprule
        \textbf{Dimension}
        & \textbf{Evaluated aspect}
        & \textbf{Subjects}
        & \textbf{Main evidence} \\
        \midrule

        Representational coverage
        & Framework capabilities expressible through the common variability model
        & Ten selected frameworks
        & Framework documentation, normalized capabilities, and GEAR features \\

        Transformation fidelity
        & Preservation of GEAR properties in framework-specific artefacts
        & Ten selected frameworks
        & Mappings, transformation reports, and generated source code \\

        Execution correctness
        & Syntactic validity, importability, and executability of generated artefacts
        & Generated implementations
        & Validation, import, and execution results \\

        Behavioral fidelity
        & Conformance of runtime behavior to the source GEAR specification
        & Generated workflow implementations
        & Expected and observed execution traces \\

        Usability and practical utility
        & Effort, understanding, and perceived usefulness
        & Future study participants
        & Task measurements, observations, and participant feedback \\

        \bottomrule
    \end{tabularx}
\end{table*}

The evaluation covers capabilities related to agent definition, configuration,
composition, orchestration, transformation, generation, and execution. It does not
assess the quality of the textual responses produced by the underlying
\gls{llm}s, model accuracy, token consumption, execution performance, or graphical
user interfaces. These aspects fall outside the current objectives, which concern
the representation, transformation, and execution of multi-agent system
configurations.

Finally, the ten frameworks considered in the representational analysis also
contributed to the construction of the common variability model. Consequently,
this part of the evaluation measures the coverage of the construction corpus. It
does not, by itself, demonstrate that GEAR generalizes to frameworks that were not
considered during the model derivation.

\subsection{Representational Coverage and Transformation Fidelity}
\label{subsec:representation-transformation}

This first evaluation axis examines two complementary properties of GEAR.
First, it determines whether the common variability model can represent the
design, configuration, and orchestration capabilities provided by the ten
selected frameworks. Second, it determines whether the semantics expressed
through GEAR are preserved when framework-specific artefacts are generated.

These two directions are evaluated using two matrices: a feature-presence
matrix for the analysis from the frameworks to GEAR, and a
semantic-preservation matrix for the analysis from GEAR to the target
frameworks. The normalized feature catalogue, the GEAR feature models, and
the framework connectors are fixed before the assessment. Consequently, this
evaluation does not modify the common model in response to the observed
results.

\paragraph{Feature-presence matrix.}

The feature-presence matrix evaluates the representational coverage of GEAR.
It answers the following question: \emph{to what extent can the capabilities
provided by each framework be represented through the common variability
model of GEAR?}

Each row corresponds to a normalized capability from the feature catalogue.
The capabilities cover agent definition, model and tool configuration, task
composition, memory, data exchange, orchestration, and execution management.
The columns correspond to GEAR and the ten selected frameworks.

Each cell is assigned one of the following values:

\begin{itemize}
    \item \textbf{F} (\emph{fully present}) indicates that the technology
    provides a construct that represents the complete semantics of the
    capability;
    
    \item \textbf{P} (\emph{partially present}) indicates that only a
    restricted form of the capability is available or that some of its
    semantics cannot be expressed; and
    
    \item \textbf{--} (\emph{absent}) indicates that no corresponding
    construct is available.
\end{itemize}

The classification of the framework cells is based on the official
documentation, public APIs, and examples associated with the reference
versions reported in Table~\ref{tab:framework-metadata}. For GEAR, the
classification is based on the features and constraints defined in
\texttt{GearAgent}, \texttt{GearModule}, and \texttt{GearWorkflow}. The
presence of a similarly named class or function is not considered sufficient
unless the expected semantics can also be represented.

Table~\ref{tab:feature-presence-matrix} presents the compact
feature-presence matrix used in the article. The detailed matrix, containing
all the individual features and variants of the common model, is provided in
the replication package.

\begin{table*}[t]
    \centering
    \caption{Feature-presence matrix used to evaluate the representational
    coverage of GEAR.}
    \label{tab:feature-presence-matrix}
    \scriptsize
    \setlength{\tabcolsep}{2pt}
    \renewcommand{\arraystretch}{1.08}
    \begin{tabularx}{\textwidth}{
        @{}
        >{\raggedright\arraybackslash}X
        ccccccccccc
        @{}
    }
        \toprule
        \textbf{Normalized capability}
        & \textbf{GEAR}
        & \textbf{CrewAI}
        & \textbf{Google ADK}
        & \textbf{LangGraph}
        & \textbf{OpenAI SDK}
        & \textbf{Microsoft AF}
        & \textbf{Strands}
        & \textbf{PydanticAI}
        & \textbf{AutoGen}
        & \textbf{Semantic K.}
        & \textbf{Haystack} \\
        \midrule
        
        Agent instructions
        & ? & ? & ? & ? & ? & ? & ? & ? & ? & ? & ? \\

        Model configuration
        & ? & ? & ? & ? & ? & ? & ? & ? & ? & ? & ? \\

        Task definition
        & ? & ? & ? & ? & ? & ? & ? & ? & ? & ? & ? \\

        Task dependencies
        & ? & ? & ? & ? & ? & ? & ? & ? & ? & ? & ? \\

        Tool association
        & ? & ? & ? & ? & ? & ? & ? & ? & ? & ? & ? \\

        Individual memory
        & ? & ? & ? & ? & ? & ? & ? & ? & ? & ? & ? \\

        Shared memory
        & ? & ? & ? & ? & ? & ? & ? & ? & ? & ? & ? \\

        Structured outputs
        & ? & ? & ? & ? & ? & ? & ? & ? & ? & ? & ? \\

        Agent delegation
        & ? & ? & ? & ? & ? & ? & ? & ? & ? & ? & ? \\

        Sequential execution
        & ? & ? & ? & ? & ? & ? & ? & ? & ? & ? & ? \\

        Parallel execution
        & ? & ? & ? & ? & ? & ? & ? & ? & ? & ? & ? \\

        Fan-out/fan-in
        & ? & ? & ? & ? & ? & ? & ? & ? & ? & ? & ? \\

        Conditional routing
        & ? & ? & ? & ? & ? & ? & ? & ? & ? & ? & ? \\

        Bounded loops
        & ? & ? & ? & ? & ? & ? & ? & ? & ? & ? & ? \\

        Asynchronous execution
        & ? & ? & ? & ? & ? & ? & ? & ? & ? & ? & ? \\

        Execution tracing
        & ? & ? & ? & ? & ? & ? & ? & ? & ? & ? & ? \\

        \bottomrule
    \end{tabularx}

    \vspace{0.3em}
        \begin{minipage}{0.98\textwidth}
            \scriptsize
            Cell values: F = fully present; P = partially present;
            -- = absent. OpenAI SDK denotes the OpenAI Agents SDK,
            and Microsoft AF denotes the Microsoft Agent Framework.
        \end{minipage}
\end{table*}

For a framework \(j\), let \(F^{*}_{j}\) denote the set of capabilities
provided by that framework and included in the functional scope of the
evaluation. Let \(R_{G,j}\) denote the subset of these capabilities that can
be represented, either fully or partially, by GEAR. The representational
coverage of GEAR for framework \(j\) is defined as:

\begin{equation}
    C_{\mathrm{repr}}(j)
    =
    \frac{|R_{G,j}|}{|F^{*}_{j}|}.
    \label{eq:representational-coverage}
\end{equation}

In addition to this overall value, the numbers of fully and partially
represented capabilities are reported separately. This distinction prevents
a partial representation from being interpreted as equivalent to a complete
one. The measure is calculated for each of the ten frameworks and aggregated
over the complete catalogue.

\paragraph{Semantic-preservation matrix.}

The semantic-preservation matrix evaluates the inverse direction, from GEAR
to the target frameworks. It answers the following question: \emph{to what
extent are the features and constraints expressed in a GEAR specification
preserved in the generated framework-specific artefacts?}

Each row corresponds to a configurable GEAR capability or constraint. When
the same capability has variants with different semantics, such as sequential
and parallel execution, these variants are represented by separate entries.
Each column corresponds to one of the ten target frameworks.

The same predefined GEAR specifications are transformed for the different
targets. No framework-specific modification is introduced into the source
configurations. This makes it possible to evaluate whether portability is
provided by the GEAR connectors rather than by manually adapting the source
specification for each framework.

Each cell of the semantic-preservation matrix is assigned one of the
following values:

\begin{itemize}
    \item \textbf{N} (\emph{native}) indicates that the target framework
    directly provides a construct that preserves the complete semantics of
    the GEAR capability;

    \item \textbf{S} (\emph{synthesized}) indicates that the connector
    generates additional code to reproduce the complete semantics of the
    capability;

    \item \textbf{A} (\emph{adapted}) indicates that the generated artefact
    supports the capability with a documented restriction or semantic
    difference;

    \item \textbf{U} (\emph{unsupported}) indicates that the connector cannot
    produce an implementation preserving the intended semantics; and

    \item \textbf{NA} (\emph{not applicable}) indicates that the capability
    is outside the execution model or functional scope of the target.
\end{itemize}

Native and synthesized transformations are both considered fully preserving.
The distinction describes the implementation mechanism, not its fidelity.
Conversely, an adapted transformation remains distinguishable because an
observable restriction or semantic difference is introduced.

Table~\ref{tab:semantic-preservation-matrix} presents the compact
semantic-preservation matrix. The complete matrix in the replication package
records the individual GEAR features, their variants, the corresponding
target constructs, and the evidence supporting each classification.

\begin{table*}[t]
    \centering
    \caption{Semantic-preservation matrix used to evaluate the transformation
    of GEAR features into framework-specific artefacts.}
    \label{tab:semantic-preservation-matrix}
    \scriptsize
    \setlength{\tabcolsep}{2pt}
    \renewcommand{\arraystretch}{1.08}
    \begin{tabularx}{\textwidth}{
        @{}
        >{\raggedright\arraybackslash}X
        cccccccccc
        @{}
    }
        \toprule
        \textbf{GEAR capability}
        & \textbf{CrewAI}
        & \textbf{Google ADK}
        & \textbf{LangGraph}
        & \textbf{OpenAI SDK}
        & \textbf{Microsoft AF}
        & \textbf{Strands}
        & \textbf{PydanticAI}
        & \textbf{AutoGen}
        & \textbf{Semantic Kernel}
        & \textbf{Haystack} \\
        \midrule

        Agent instructions
        & ? & ? & ? & ? & ? & ? & ? & ? & ? & ? \\

        Model configuration
        & ? & ? & ? & ? & ? & ? & ? & ? & ? & ? \\

        Task definition
        & ? & ? & ? & ? & ? & ? & ? & ? & ? & ? \\

        Task dependencies
        & ? & ? & ? & ? & ? & ? & ? & ? & ? & ? \\

        Tool association
        & ? & ? & ? & ? & ? & ? & ? & ? & ? & ? \\

        Individual memory
        & ? & ? & ? & ? & ? & ? & ? & ? & ? & ? \\

        Shared memory
        & ? & ? & ? & ? & ? & ? & ? & ? & ? & ? \\

        Structured outputs
        & ? & ? & ? & ? & ? & ? & ? & ? & ? & ? \\

        Agent delegation
        & ? & ? & ? & ? & ? & ? & ? & ? & ? & ? \\

        Sequential execution
        & ? & ? & ? & ? & ? & ? & ? & ? & ? & ? \\

        Parallel execution
        & ? & ? & ? & ? & ? & ? & ? & ? & ? & ? \\

        Fan-out/fan-in
        & ? & ? & ? & ? & ? & ? & ? & ? & ? & ? \\

        Conditional routing
        & ? & ? & ? & ? & ? & ? & ? & ? & ? & ? \\

        Bounded loops
        & ? & ? & ? & ? & ? & ? & ? & ? & ? & ? \\

        Natural-language stop condition
        & ? & ? & ? & ? & ? & ? & ? & ? & ? & ? \\

        Asynchronous execution
        & ? & ? & ? & ? & ? & ? & ? & ? & ? & ? \\

        Execution tracing
        & ? & ? & ? & ? & ? & ? & ? & ? & ? & ? \\

        \bottomrule
    \end{tabularx}

    \vspace{0.3em}
        \begin{minipage}{0.98\textwidth}
            \scriptsize
            Cell values: N = native; S = synthesized; A = adapted;
            U = unsupported; NA = not applicable. OpenAI SDK denotes
            the OpenAI Agents SDK, and Microsoft AF denotes the
            Microsoft Agent Framework.
        \end{minipage}
\end{table*}

For a target framework \(j\), let \(n_{N,j}\), \(n_{S,j}\),
\(n_{A,j}\), and \(n_{U,j}\) denote the numbers of native, synthesized,
adapted, and unsupported transformations, respectively. Entries classified
as not applicable are excluded from the calculations. The number of
applicable transformations is therefore:

\begin{equation}
    n_{\mathrm{app},j}
    =
    n_{N,j}
    +
    n_{S,j}
    +
    n_{A,j}
    +
    n_{U,j}.
    \label{eq:applicable-transformations}
\end{equation}

The semantic-preservation coverage measures the proportion of GEAR
capabilities whose complete semantics are preserved, either through a native
construct or through synthesized code:

\begin{equation}
    C_{\mathrm{pres}}(j)
    =
    \frac{n_{N,j}+n_{S,j}}
         {n_{\mathrm{app},j}}.
    \label{eq:preservation-coverage}
\end{equation}

Adapted capabilities are not included in this measure because their
transformation introduces a documented semantic difference. They are
reported separately through the adapted coverage:

\begin{equation}
    C_{\mathrm{adapt}}(j)
    =
    \frac{n_{A,j}}
         {n_{\mathrm{app},j}}.
    \label{eq:adapted-coverage}
\end{equation}

The unsupported coverage measures the proportion of applicable GEAR
capabilities for which no valid transformation is available:

\begin{equation}
    C_{\mathrm{unsupported}}(j)
    =
    \frac{n_{U,j}}
         {n_{\mathrm{app},j}}.
    \label{eq:unsupported-coverage}
\end{equation}

These three measures provide a complete decomposition of the applicable
transformations:

\begin{equation}
    C_{\mathrm{pres}}(j)
    +
    C_{\mathrm{adapt}}(j)
    +
    C_{\mathrm{unsupported}}(j)
    =
    1.
    \label{eq:transformation-decomposition}
\end{equation}

Because a transformation limitation must not remain silent, the assessment
also examines whether GEAR reports an explicit diagnostic for adapted and
unsupported capabilities. Let \(n_{\mathrm{diag},j}\) denote the number of
adapted or unsupported transformations for which the connector produces an
explicit diagnostic. Diagnostic coverage is defined as:

\begin{equation}
    C_{\mathrm{diag}}(j)
    =
    \frac{n_{\mathrm{diag},j}}
         {n_{A,j}+n_{U,j}}.
    \label{eq:diagnostic-coverage}
\end{equation}

When no adapted or unsupported capability is observed for a framework,
diagnostic coverage is reported as not applicable. The measures are reported
separately for each framework. Macro-aggregated results give the same weight
to each framework, whereas micro-aggregated results are computed from all
applicable matrix entries.

\paragraph{Evidence and traceability.}

Each classification is associated with evidence identifying the normalized
capability or GEAR feature, the corresponding framework construct, the
applicable documentation or API entry, the connector mapping, and the
generated source-code fragment. This information is stored in the
replication package together with the complete matrices and the scripts used
to calculate the reported measures.

The two matrices provide complementary evidence. The feature-presence matrix
determines whether the abstraction offered by GEAR covers the capabilities
available in the selected frameworks. The semantic-preservation matrix
determines how these abstractions are translated back into concrete
framework-specific implementations.

This analysis remains static: it establishes whether the required
information and semantics are represented in the common model and in the
generated source code. The following subsection evaluates whether the
generated artefacts are syntactically valid and executable, and whether
transformations classified as native, synthesized, or adapted are confirmed
by their observed runtime behavior.

\subsection{Execution Correctness and Behavioral Fidelity}
\label{subsec:execution-behavior}

This evaluation axis determines whether the framework-specific artefacts generated
by GEAR can be executed and whether their runtime behavior conforms to the source
GEAR specifications. While the previous analysis examines mappings and generated
code statically, this analysis provides dynamic evidence through validation,
execution, and trace comparison.

\paragraph{Execution correctness.}
Each evaluation specification is transformed into an implementation for every
applicable target framework. The generated artefacts are then assessed through
four successive stages:

\begin{enumerate}
    \item \emph{syntactic validation}, which checks that the generated source files
    can be parsed or compiled;
    \item \emph{dependency and import validation}, which checks that the generated
    project can load the required framework components;
    \item \emph{initialization validation}, which checks that the agents, tasks,
    modules, and workflows can be instantiated; and
    \item \emph{execution validation}, which checks that the generated workflow can
    run to completion without an error attributable to the transformation.
\end{enumerate}

A generated implementation is considered executable only if it passes all four
stages. Failures are recorded with their stage, framework, specification, and
error message. When a GEAR property is unsupported by a target framework, the
expected behavior is an explicit diagnostic during validation or generation,
rather than the production of a silently incomplete implementation.

For a target framework \(t\), execution correctness is measured as:

\begin{equation}
    \mathrm{ExecRate}_t =
    \frac{N^{\mathrm{executed}}_t}
         {N^{\mathrm{applicable}}_t},
\end{equation}

where \(N^{\mathrm{applicable}}_t\) is the number of implementations expected to
be generated and executed for framework \(t\), and
\(N^{\mathrm{executed}}_t\) is the number that successfully pass all validation
and execution stages. Stage-specific success rates are also reported to
distinguish generation, import, initialization, and runtime failures.

\paragraph{Behavioral fidelity.}
Successful execution alone does not demonstrate that the generated implementation
preserves the workflow semantics. Behavioral fidelity is therefore evaluated by
comparing the observed execution trace with an expected trace derived from the
source GEAR specification.

The evaluation scenarios are constructed to exercise the orchestration constructs
defined by \texttt{GearWorkflow}, including sequential execution, dependencies,
parallel branches, synchronization points, conditional transitions, and bounded
loops. Where applicable, scenarios also exercise agent delegation and the
transmission of outputs between dependent tasks.

For each scenario, the source specification is converted into a set of observable
constraints. These constraints define, for example:

\begin{itemize}
    \item which tasks must be executed;
    \item which task must precede another;
    \item which tasks may execute concurrently;
    \item which branches may be activated under a given condition;
    \item where parallel branches must be synchronized;
    \item how many times a bounded loop must be executed; and
    \item which outputs must be made available to dependent tasks.
\end{itemize}

The generated implementations are instrumented using a common trace format. Each
event records at least the target framework, scenario, agent, task, event type,
and execution order. Framework-specific traces are normalized before comparison
so that differences in logging formats do not affect the assessment.

Controlled task outputs and deterministic conditions are used during execution.
This isolates orchestration behavior from the variability of the underlying
\gls{llm}s. Consequently, the evaluation concerns the behavior produced by the
generated orchestration logic and not the quality of the textual responses
generated by a model.

Let \(K_s\) denote the set of behavioral constraints derived from specification
\(s\), and let \(K^{\mathrm{sat}}_{s,t}\) denote the constraints satisfied by the
trace observed for target framework \(t\). Behavioral fidelity is measured as:

\begin{equation}
    \mathrm{TraceConformance}_{s,t} =
    \frac{|K^{\mathrm{sat}}_{s,t}|}{|K_s|}.
\end{equation}

A trace is considered fully conformant when all applicable constraints are
satisfied. Violations are reported individually and classified according to the
affected property, such as ordering, dependency, branching, synchronization,
iteration, or data transmission.

The same source specifications, controlled outputs, and trace constraints are
used across the target frameworks. This enables a cross-framework comparison
without introducing framework-specific changes into the GEAR configurations.
The execution results are also compared with the semantic preservation matrix.
A property classified as \emph{direct} or \emph{equivalent} is expected to
produce a conformant trace, whereas a \emph{partial} classification must exhibit
only the limitations declared during the static analysis. Any additional runtime
deviation is recorded as a discrepancy between the static mapping and the
observed behavior.